\textbf{Phần 1. Trắc nghiệm lựa chọn}
\setcounterpageref{ex}{0}

\begin{ex} %[4D5N23]
Kết quả của phép trừ $678\,950 - 234\,510$ là
\choice
{$444\,430$}
{$444\,540$}
{$434\,440$}
{\True $444\,440$}
\loigiai{
$678\,950 - 234\,510 = 444\,440$.\\
\textbf{Đáp án của câu hỏi là D}
}
\end{ex}

\begin{ex} %[4D5N23]
Kết quả của phép tính $5\,000\,000 - 2\,000\,000$ là
\choice
{$2\,000\,000$}
{$4\,000\,000$}
{\True $3\,000\,000$}
{$3\,500\,000$}
\loigiai{
$5\,000\,000 - 2\,000\,000 = 3\,000\,000$.\\
\textbf{Đáp án của câu hỏi là C}
}
\end{ex}

\begin{ex} %[4D5H23]
Kết quả của phép trừ $10\,000\,000 - 3\,450\,000$ là
\choice
{$6\,450\,000$}
{$6\,650\,000$}
{$7\,550\,000$}
{\True $6\,550\,000$}
\loigiai{
$10\,000\,000 - 3\,450\,000 = 6\,550\,000$.\\
\textbf{Đáp án của câu hỏi là D}
}
\end{ex}

\begin{ex} %[4D5H23]
Tìm $x$ biết: $750\,000 - x = 280\,000$.
\choice
{\True $470\,000$}
{$460\,000$}
{$480\,000$}
{$570\,000$}
\loigiai{
$x = 750\,000 - 280\,000 = 470\,000$.\\
\textbf{Đáp án của câu hỏi là A}
}
\end{ex}

\begin{ex} %[4D5H23]
Một kho chứa $45\,000\text{ kg}$ thóc. Sau khi xuất đi $18\,500\text{ kg}$ thóc thì trong kho còn lại bao nhiêu ki-lô-gam thóc?
\choice
{$25\,500\text{ kg}$}
{\True $26\,500\text{ kg}$}
{$27\,500\text{ kg}$}
{$26\,000\text{ kg}$}
\loigiai{
$45\,000 - 18\,500 = 26\,500\text{ kg}$.\\
\textbf{Đáp án của câu hỏi là B}
}
\end{ex}

\begin{ex} %[4D5H23]
Hiệu của số bé nhất có 7 chữ số và số lớn nhất có 5 chữ số là
\choice
{\True $900\,001$}
{$900\,001$}
{$900\,000$}
{$899\,999$}
\loigiai{
$1\,000\,000 - 99\,999 = 900\,001$.\\
\textbf{Đáp án của câu hỏi là A}
}
\end{ex}

\begin{ex} %[4D5V23]
Một tài khoản ngân hàng có $85\,000\,000$ đồng. Chủ tài khoản rút ra $25\,000\,000$ đồng để chi tiêu và chuyển khoản $30\,000\,000$ đồng để thanh toán tiền hàng. Số tiền còn lại trong tài khoản là
\choice
{$25\,000\,000$ đồng}
{\True $30\,000\,000$ đồng}
{$35\,000\,000$ đồng}
{$40\,000\,000$ đồng}
\loigiai{
Còn lại: $85\,000\,000 - (25\,000\,000 + 30\,000\,000) = 30\,000\,000$ đồng.\\
\textbf{Đáp án của câu hỏi là B}
}
\end{ex}

\begin{ex} %[4D5V23]
Tính giá trị của biểu thức: $120\,000\,000 - (45\,000\,000 + 35\,000\,000)$.
\choice
{$30\,000\,000$}
{$50\,000\,000$}
{\True $40\,000\,000$}
{$45\,000\,000$}
\loigiai{
$120\,000\,000 - 80\,000\,000 = 40\,000\,000$.\\
\textbf{Đáp án của câu hỏi là C}
}
\end{ex}

\begin{ex} %[4D5V23]
Dân số thành phố A là $3\,450\,000$ người, dân số thành phố B là $2\,150\,000$ người. Dân số thành phố A nhiều hơn thành phố B bao nhiêu người?
\choice
{$1\,200\,000$ người}
{$1\,400\,000$ người}
{$1\,500\,000$ người}
{\True $1\,300\,000$ người}
\loigiai{
$3\,450\,000 - 2\,150\,000 = 1\,300\,000$ người.\\
\textbf{Đáp án của câu hỏi là D}
}
\end{ex}

\begin{ex} %[4D5V23]
Một xí nghiệp may theo kế hoạch phải may $150\,000$ bộ quần áo. Đợt 1 may được $45\,000$ bộ, đợt 2 may được $55\,000$ bộ. Để hoàn thành kế hoạch xí nghiệp còn phải may số bộ quần áo là
\choice
{\True $50\,000$ bộ}
{$40\,000$ bộ}
{$60\,000$ bộ}
{$55\,000$ bộ}
\loigiai{
Còn lại: $150\,000 - (45\,000 + 55\,000) = 50\,000$ bộ.\\
\textbf{Đáp án của câu hỏi là A}
}
\end{ex}

\begin{ex} %[4D5V23]
Tìm số bị trừ biết số trừ là $3\,450\,000$ và hiệu là $2\,550\,000$.
\choice
{$5\,900\,000$}
{$6\,100\,000$}
{\True $6\,000\,000$}
{$5\,000\,000$}
\loigiai{
Số bị trừ: $3\,450\,000 + 2\,550\,000 = 6\,000\,000$.\\
\textbf{Đáp án của câu hỏi là C}
}
\end{ex}

\begin{ex} %[4D5C23]
Khi thực hiện phép trừ một số có 6 chữ số cho một số có 5 chữ số, một bạn đã viết nhầm số trừ lùi sang trái 1 cột nên được hiệu là $450\,000$. Biết số bị trừ là $650\,000$. Tìm số trừ đúng.
\choice
{$10\,000$}
{\True $20\,000$}
{$30\,000$}
{$40\,000$}
\loigiai{
Viết lùi sang trái 1 cột tức là gấp 10 lần số trừ.\\
Số trừ sai là: $650\,000 - 450\,000 = 200\,000$.\\
Số trừ đúng là: $200\,000 : 10 = 20\,000$.\\
\textbf{Đáp án của câu hỏi là B}
}
\end{ex}

\textbf{Phần 2. Trắc nghiệm đúng sai}
\setcounterpageref{ex}{0}

\begin{ex} %[4D5H23]
Xét tính đúng sai của các phép trừ sau:
\choiceTF
{\True $1\,000\,000 - 400\,000 = 600\,000$.}
{$850\,000 - 320\,000 = 540\,000$.}
{\True $850\,000 - 320\,000 = 530\,000$.}
{\True $45\,678\,900 - 15\,678\,900 = 30\,000\,000$.}
\loigiai{
\begin{itemize}
    \item $1\,000\,000 - 400\,000 = 600\,000$. Đúng.
    \item $850\,000 - 320\,000 = 530\,000 \neq 540\,000$. Ý 2 Sai, Ý 3 Đúng.
    \item $45\,678\,900 - 15\,678\,900 = 30\,000\,000$. Đúng.
\end{itemize}
\begin{itemize}
    \item a) \textbf{Đúng}.
    \item b) \textbf{Sai}.
    \item c) \textbf{Đúng}.
    \item d) \textbf{Đúng}.
\end{itemize}
}
\end{ex}

\begin{ex} %[4D5V23]
Một xí nghiệp có $120\,000$ tấn than trong kho: Ngày thứ nhất xuất $35\,000$ tấn than, ngày thứ hai xuất $45\,000$ tấn than. Xét tính đúng sai:
\choiceTF
{\True Tổng số than xuất trong 2 ngày là $80\,000$ tấn.}
{\True Số than xuất ngày thứ hai nhiều hơn ngày thứ nhất $10\,000$ tấn.}
{Sau hai ngày kho còn lại $50\,000$ tấn than.}
{\True Sau hai ngày kho còn lại $40\,000$ tấn than.}
\loigiai{
\begin{itemize}
    \item Đã xuất: $35\,000 + 45\,000 = 80\,000$ tấn. Đúng.
    \item Hiệu: $45\,000 - 35\,000 = 10\,000$ tấn. Đúng.
    \item Còn lại: $120\,000 - 80\,000 = 40\,000$ tấn $\neq 50\,000$ tấn. Ý 3 Sai, Ý 4 Đúng.
\end{itemize}
\begin{itemize}
    \item a) \textbf{Đúng}.
    \item b) \textbf{Đúng}.
    \item c) \textbf{Sai}.
    \item d) \textbf{Đúng}.
\end{itemize}
}
\end{ex}

\textbf{Phần 3. Trả lời ngắn}
\setcounterpageref{ex}{0}

\begin{ex} %[4D5H23]
Tính hiệu: $900\,000 - 350\,000$.
\shortans{550000}
\loigiai{
\textbf{Đáp số:} $550000$\\
$550\,000$.
}
\end{ex}

\begin{ex} %[4D5H23]
Tính hiệu: $50\,000\,000 - 18\,000\,000$.
\shortans{32000000}
\loigiai{
\textbf{Đáp số:} $32000000$\\
$32\,000\,000$.
}
\end{ex}

\begin{ex} %[4D5V23]
Tìm $x$ biết: $1\,500\,000 - x = 850\,000$.
\shortans{650000}
\loigiai{
\textbf{Đáp số:} $650000$\\
$x = 1\,500\,000 - 850\,000 = 650\,000$.
}
\end{ex}

\begin{ex} %[4D5V23]
Tính giá trị biểu thức: $500\,000 - 120\,000 - 180\,000$.
\shortans{200000}
\loigiai{
\textbf{Đáp số:} $200000$\\
$500\,000 - 300\,000 = 200\,000$.
}
\end{ex}

\begin{ex} %[4D5V23]
Một công ty dự định thu về $500$ triệu đồng, thực tế chỉ thu được $420$ triệu đồng. Hỏi công ty thu thiếu bao nhiêu triệu đồng? (Chỉ ghi số triệu).
\shortans{80}
\loigiai{
\textbf{Đáp số:} $80$\\
$500 - 420 = 80$ triệu.
}
\end{ex}

\begin{ex} %[4D5C23]
Hiệu của hai số là $45\,000$. Nếu thêm vào số bị trừ $5000$ đơn vị và bớt ở số trừ $5000$ đơn vị thì hiệu mới là bao nhiêu?
\shortans{55000}
\loigiai{
\textbf{Đáp số:} $55000$\\
$45\,000 + 5000 + 5000 = 55\,000$.
}
\end{ex}
